\section{Conditions of Corporeal Manifestation}

In 1912, when Guenon was in his mid-twenties, he published a densely argued article on the conditions of corporeal manifestation. It abruptly ends; apparently a Part 2 was planned to conclude the essay, but never appeared. Rather than try to extrapolate the second half (which would be difficult without knowing his source material), I'd rather use what is available and then apply it in a way to understand the natural or human order.

Guenon is one-pointedly focused on the cosmic or metaphysical order; when he does talk about the human order, it is usually to criticize various aspects of it. Nevertheless, beyond the Brahmanic type of spirituality that he undoubtedly represents, there are many more whose earthly task brings them into the areas of administration or economic activity. There certainly is a spirituality appropriate for them. Guenon does indeed point out that the human order and the cosmic order are not really separated, since they continuously react on each other (``Solidification of the World'', Reign of Quantity).

Guenon writes that there are five conditions which corporeal existence is subject to:

\begin{itemize}
\item space
\item time
\item matter
\item form
\item life
\end{itemize}

He summarizes: ``A body is a material form living in time and space.'' Space and time are clearly necessary for there to be any consciousness of material objects. Neither space nor time is discovered through observation, but rather they are necessary conditions for there to be observations at all.

By including matter and form, Guenon reaffirms the hylomorphism of the Western Tradition. Taking some cues from Guenon and applying them to Plato's divided line, understood ontologically rather than epistemologically, we can attempt to understand the conditions of material existence. Plato makes a distinction between philosophical ideas and mathematical ideas both because of their forms of knowledge and that the latter are closer to the material world. Guenon claims the Point, without dimension, generates all the other dimensions. So we already start with a mathematicization of the world.

The two branches of maths, arithmetic and geometry, then correspond to time and space, respectively. Since time is more subtle (for example, a thought exists in time but not in space), it is ontologically prior to space. These are the conditions for the appearance of things. The actuality of things is then based on form and matter, which we can consider as two opposing poles.

For God, essence and existence are identical, but not so for the world. The Absolute is the Point, and the Infinite relates to all the possibilities. In separating those two aspects, there is a clearing for the world to appear. Every corporeal being partakes, then, in both form and matter.


\flrightit{Posted on 2012-09-27 by Cologero}

